%\usepackage[latin1]{inputenc}

\usepackage{graphicx}
\usepackage{rotating}
%\setbeamertemplate{caption}[numbered]
\usepackage{hyperref}
\usepackage{caption}
\usepackage[normalem]{ulem}
%\mode<presentation>
\usepackage{wasysym}
\usepackage{amsmath}
\usepackage{mathtools}
\usepackage[skins,theorems]{tcolorbox}
% \usepackage[,dvipsnames]{xcolor}
\tcbset{highlight math style={enhanced,
  colframe=red,colback=white,arc=0pt,boxrule=1pt}}

\newcommand{\normal}[2]{\ensuremath{\mathcal{N}\left (#1,#2 \right )}}
\newcommand{\Gaussian}[3]{\ensuremath{\frac{1}{\sqrt{2\pi}#3}
\text{exp}\left \{-\frac{1}{2#3^2} (#1-#2)^2 \right \}}}
\newcommand{\argmax}{\operatornamewithlimits{argmax}}
\newcommand{\expo}[1]{\ensuremath{\text{exp}\left \{ #1 \right \}}}
\newcommand{\studentt}[4]{\ensuremath{\mathcal{T}_{#4}(#1,#2,#3)}}
\newcommand{\vv}[1]{\boldsymbol{#1}}
\newcommand{\Prb}{\ensuremath{\mathbb{P}}}
\newcommand{\studenttk}[4]{\ensuremath{\mathcal{T}_{#4}(#1,#2,#3)}}
\newcommand{\argmin}{\operatornamewithlimits{argmin}}
\newcommand{\NIG}{\mathcal{NIG}}
\newcommand{\N}{\mathcal{N}}
\newcommand{\T}{\mathcal{T}}
\newcommand{\IG}{\mathcal{IG}}
\newcommand{\IW}{\mathcal{IW}}
\newcommand{\NNIW}{\mathcal{NNIW}}
\newcommand{\vct}{\text{vec}}
\newcommand{\NN}{\mathcal{NN}}
\newcommand{\tr}{\text{tr}}
\newcommand{\di}[2]{\ensuremath{ #1^{(#2)}}}
\newcommand{\dd}[1]{\ensuremath{ #1^{(i)}}}
\newcommand{\Di}[2]{\ensuremath{ \vv{#1}^{(#2)}}}

\newcommand{\E}[1]{\ensuremath{\mathbb{E}[#1]}}
\newcommand{\Var}[1]{\mathrm{Var}[#1]}
\newcommand{\Cov}[2]{\mathrm{Cov}[#1,#2]}
\newcommand{\Cor}[2]{\mathrm{Cor}[#1,#2]}

\usepackage[]{algorithm2e}
% \setbeamertemplate{navigation symbols}{}
\usepackage{tcolorbox}


%\titlegraphic{\includegraphics[width=0.3\paperwidth]{\string~/Dropbox/teaching/clemson-academic.png}}
\titlegraphic{\includegraphics[width=0.1\paperwidth]{./figs/crest.pdf}}
\setbeamertemplate{title page}[empty]

\setbeamerfont{subtitle}{size=\small}

\setbeamercovered{transparent}

\definecolor{clemsonpurple}{HTML}{522D80}
\definecolor{stablue}{HTML}{0052cc}
\definecolor{stared}{HTML}{ff4d4d}
\definecolor{clemsonorange}{HTML}{F66733}
\newcommand{\empha}[1]{\textbf{\textcolor{stablue}{#1}}}
\setbeamercolor{frametitle}{fg=stablue,bg=white}
\setbeamercolor{title}{fg=stablue,bg=white}
\setbeamercolor{local structure}{fg=stablue}
\setbeamercolor{section in toc}{fg=stablue,bg=white}
\setbeamercolor{subsection in toc}{fg=stared,bg=white}
\setbeamercolor{item projected}{fg=stablue,bg=white}
\setbeamertemplate{itemize item}{\color{stablue}$\bullet$}
\setbeamertemplate{itemize subitem}{\color{stablue}\scriptsize{$\bullet$}}
\let\Tiny=\tiny

%\makeatletter
%\setbeamertemplate{footline}{%
%\leavevmode
%\vbox{\begin{beamercolorbox}[dp=1.25ex,ht=2.75ex]{fg=black}%
%  \hspace*{1em}\insertsectionhead%
%  \ifx\insertsubsectionhead\@empty\relax\else$\quad\mid\quad$\insertsubsectionhead\fi :
%  \end{beamercolorbox}%
%  }%
%}
%\makeatother

\setbeamercolor{footercl}{fg=white,bg=stablue}
\setbeamerfont{stafont}{size = \large}
\setbeamerfont{footerfont}{size = \tiny}
% \makeatother
% \setbeamertemplate{footline}
% {
%   \leavevmode%
%   \hbox{%
%   \begin{beamercolorbox}[wd=.5\paperwidth,ht=5ex,dp=1ex, left]{footercl}%
%     \usebeamerfont*{stafont}\hspace*{1ex}\insertshortinstitute
%   \end{beamercolorbox}%
%   \begin{beamercolorbox}[wd=.25\paperwidth,ht=5ex,dp=1ex,center]{footercl}%
%     \usebeamerfont*{footerfont}\insertshorttitle\hspace*{1ex} \insertframenumber{}
%   \end{beamercolorbox}%
%   \begin{beamercolorbox}[wd=.25\paperwidth,ht=5ex,dp=1ex,right]{footercl}%
%    \includegraphics[height=5ex]{stalogo.png}
%   \end{beamercolorbox}}%
%   \vskip0pt%
% }
% \makeatletter
\makeatletter
\setbeamertemplate{footline}
{
  \leavevmode%
  \hbox{%
  \fontsize{13}{13}\fontfamily{ppl}\selectfont
  \begin{beamercolorbox}[wd=.5\paperwidth,ht=2.25ex,dp=1ex,left]{footercl}%
    \usebeamerfont{author in head/foot}\hspace*{1ex}\insertshortinstitute
   \end{beamercolorbox}%
   \begin{beamercolorbox}[wd=.25\paperwidth,ht=2.25ex,dp=1ex,right]{footercl}%
    % \hfill\hfill\hfill\hfill\hfill\hfill\hfill\hfill\hfill\hfill
    \parbox{.25\paperwidth}{\fontfamily{cmss}\selectfont{\hfill\hfill \usebeamerfont{footerfont}\insertshorttitle~\insertframenumber{}}}
  \end{beamercolorbox}%
  \begin{beamercolorbox}[wd=.25\paperwidth,ht=2.25ex,dp=1ex,left]{footercl}%
    \parbox{.25\paperwidth}{\hfill\includegraphics[height=1cm]{./figs/stalogo.png}}% original: 2ex
  \end{beamercolorbox}}%
  \vskip0pt%
}
\makeatother

\setbeamertemplate{navigation symbols}{}

% \AtBeginDocument{\author[L. Fang]{Lei Fang} \institute[www.st-andrews.ac.uk]{School of Computer Science, University of St Andrews}}

% \newcommand{\Ffootline}{%                   %%defines a new command called \Ffootline
% %\insertshortauthor,                         %%puts the abbreviated form of the author's name in the left corner
% \insertshorttitle,
% \insertshortinstitute, 
% \insertshortdate %%puts the abbreviated form of the author's institution in the middle
% \hfill
% \insertsection,
% \insertframenumber/\inserttotalframenumber} %%includes the current slide number over the total slide number in the right corner
% \setbeamertemplate{footline}{%              %%sets the options for the footline
% \usebeamerfont{structure}                   %%uses the same fonts adopted for the structure of the presentation 
% \Tiny\hspace*{4mm} \Ffootline \hspace{4mm}  %%sets the size of the font to Tiny and includes the content of the \Ffootline
% }                                           %%command leaving a margin of 4mm to the right and left of the content.



\AtBeginPart{}
\AtBeginSection{}
\AtBeginSubsection{}
\AtBeginSubsubsection{}
\setlength{\emergencystretch}{0em}
\setlength{\parskip}{0pt}